% ============================================================
%  Section: Database Design and Storage
% ============================================================

\subsection{Database System}

The project uses \textbf{SQLite} as its relational database engine. SQLite is an
embedded, serverless database that stores the entire database as a single
\texttt{.db} file. This choice was made for several practical reasons:

\begin{itemize}
    \item No server installation or configuration required.
    \item Native Python support via the built-in \texttt{sqlite3} module.
    \item Sufficient performance for analytical queries on millions of rows.
    \item Portability: the database file can be shared or recreated from the pipeline.
\end{itemize}

\subsection{Schema Design}

% Add your ER diagram here once you have it.
% Placeholder:
\begin{figure}[H]
    \centering
    % \includegraphics[width=0.85\textwidth]{figures/er_diagram.png}
    \fbox{\parbox{0.85\textwidth}{\centering\vspace{2cm}
    \textit{[Insert ER diagram here]}
    \vspace{2cm}}}
    \caption{Entity-Relationship Diagram for the MTA Transit Pulse database}
    \label{fig:er_diagram}
\end{figure}

The database consists of the following tables:

\subsubsection{\texttt{stations}}
Stores one row per unique station, normalized from the raw STATION column.

\subsubsection{\texttt{turnstiles}}
Stores one row per unique turnstile unit (identified by C/A + UNIT + SCP).
Each turnstile is linked to a station via foreign key.

\subsubsection{\texttt{readings}}
Stores the individual counter readings with parsed timestamps.
This is the largest table, with one row per four-hour interval per turnstile.

\subsubsection{\texttt{daily\_ridership}}
A derived/aggregated table that pre-computes daily net entries and exits per
station. This table is rebuilt by the analysis pipeline and powers the
Streamlit dashboard queries.

\subsection{Key SQL Patterns}

% Once you've written your queries, paste 1-2 representative ones here.
% Example placeholder:

\begin{lstlisting}[style=sqlstyle, caption={Example: Top 10 stations by average daily entries}]
SELECT
    s.station_name,
    AVG(dr.net_entries) AS avg_daily_entries
FROM daily_ridership dr
JOIN stations s ON dr.station_id = s.id
GROUP BY s.station_name
ORDER BY avg_daily_entries DESC
LIMIT 10;
\end{lstlisting}
